% ID: SEL-D0076-OA-E03
% ESTADO: cotejada. Segunda pasada visual y borrador compilado revisados.
% Fuente: Selectividad/sel_2008_matematicas_aplicadas/Mat CCSS 4 - Examen y criterios de correción.pdf
% SHA-256: f1040adb7b9fb01c26219800e6212b9a26b93bb9d5e933743d18ee6ac26e163b
% Página: 1; opción A; ejercicio original 3.

\noindent\underline{Parte I}

En una población, donde el $45\%$ son hombres y el resto mujeres, se sabe que el $10\%$ de los hombres y el $8\%$ de las mujeres son inmigrantes.
\begin{enumerate}[label=\alph*)]
\item \textbf{(1 punto)} \textquestiondown{}Qué porcentaje de inmigrantes hay en esta población?
\item \textbf{(1 punto)} Si se elige, al azar, un inmigrante de esta población, \textquestiondown{}cuál es la probabilidad de que sea hombre?
\end{enumerate}
\par\noindent\begin{minipage}{\linewidth}
\noindent\underline{Parte II}

\textbf{(2 puntos)} Tomada al azar una muestra de 90 alumnos de un Instituto se encontró que un tercio habla inglés.

Halle, con un nivel de confianza del $97\%$, un intervalo de confianza para estimar la proporción de alumnos de ese Instituto que habla inglés.
\end{minipage}\par
